% Sección 1: Introducción y objetivos (en APA 7 el título repetido actúa como encabezado de nivel 1)
% evidencia:E01,E02,E03

\section{MediFlow}

La gestión operativa de un hospital de alta complejidad exige articular la anticipación de llegadas diarias con el comportamiento de la permanencia intrahospitalaria y la distribución horaria de los recursos de atención. En el Hospital Universidad del Norte (HUN), la disponibilidad de un repositorio clínico estructurado bajo el estándar Observational Medical Outcomes Partnership Common Data Model (OMOP CDM) versión 5.4 \parencite{ohdsi54} abre la posibilidad de analizar trayectorias agregadas sin exponer registros individuales de pacientes; sin embargo, dicha representación administrativa carece de tablas de capacidad física instalada, camas censales habilitadas y dotación contractual de personal por turno \parencite{mediflow2026informe}. Ante esa restricción estructural, MediFlow se concibe como una herramienta de apoyo a la planificación que sigue el hilo conductor de \textbf{comprender, predecir y decidir}: reconstruye el flujo histórico por etapas, evalúa la incertidumbre de la demanda con validación temporal estricta y traduce pronósticos y reglas de permanencia en escenarios paramétricos interpretables con supuestos explícitos.

El aporte defendible del proyecto reside en la integración trazable de cuatro capas técnicas que permanecen desacopladas de los datos individuales en el despliegue: extracción SQL agregada de solo lectura en origen, modelado predictivo con evaluación temporal de origen móvil y explicabilidad aditiva, análisis operacional basado en teoría de colas y balance de flujo, e interfaces interactivas (aplicación web FastAPI/React y tablero histórico en Power BI). En coherencia con la evidencia disponible, el sistema no afirma haber demostrado una reducción clínica observada en los tiempos de espera ni una asignación óptima sobre recursos reales del hospital; todo escenario cuantificado se presenta como hipótesis operacional sujeta a contraste con la capacidad y las marcas temporales del sistema de información hospitalario (HIS).

\subsection{Objetivos y alcance editorial}

Como síntesis editorial del trabajo desarrollado por el equipo para el Data Challenge HUN 2026, el objetivo general consiste en analizar el flujo de pacientes y apoyar la planificación operacional del HUN mediante pronósticos de demanda con intervalos de incertidumbre y simulaciones interpretables bajo supuestos explícitos. De este propósito se desprenden cinco objetivos específicos: (a) reconstruir la demanda por servicio, las etapas de atención, las transiciones entre áreas y el censo horario a partir de agregados OMOP; (b) auditar la calidad, la completitud, la consistencia temporal y el aporte empírico de factores externos mediante protocolos anti-fuga; (c) contrastar modelos predictivos de corto y largo horizonte con métricas libres de explosión por ceros y evaluación temporal estricta; (d) cuantificar escenarios paramétricos de dotación equivalente, permanencia promedio y redistribución intradía de egresos; y (e) comunicar las capacidades implementadas, las diferencias entre versiones experimentales y los límites que condicionan una eventual validación clínica. Esta formulación ordena la exposición técnica del informe sin introducir requisitos ajenos a las bases oficiales del certamen.

\subsection{Correspondencia con los requisitos oficiales del certamen}

Las exigencias normativas del proyecto provienen exclusivamente de los términos y condiciones de participación y del enunciado técnico oficial del Data Challenge HUN 2026 \parencite{hun2026terminos,hun2026reto}. Los documentos oficiales establecen como entregables la documentación técnica que consolide metodología, análisis, resultados, conclusiones y recomendaciones; un repositorio en GitHub con el código y recursos autorizados; un enlace interactivo en Power BI; y la sustentación oral cuando corresponda. El enunciado permite abordar uno o varios componentes del flujo hospitalario y no exige emplear todas las tablas clínicas como variables predictoras ni desplegar un piloto asistencial previo en producción.

\begin{table}[htbp]
\caption{Criterios oficiales de evaluación y desempates del Data Challenge HUN 2026}
\label{tab:criterios_oficiales}
{\singlespacing\small
\begin{tabularx}{\textwidth}{@{}L{0.46\textwidth} r Y@{}}
\toprule
Criterio oficial de evaluación & Peso & Bloque principal del informe \\
\midrule
Comprensión del problema y análisis del flujo hospitalario & 15\,\% & Introducción; Datos y reconstrucción del flujo \\
Calidad, preparación e integración de los datos & 15\,\% & Datos y flujo; Calidad, preparación y externos \\
Modelo analítico y/o predictivo y rigor metodológico & 20\,\% & Método analítico y validación; Resultados predictivos \\
Optimización y propuestas de mejora operacional & 20\,\% & Optimización operacional y simulaciones paramétricas \\
Integración de la solución, visualización y dashboard & 15\,\% & Integración, aplicación web y visualización \\
Ingeniería de software, reproducibilidad y documentación & 5\,\% & Ingeniería de software, reproducibilidad y privacidad \\
Presentación, interpretación y recomendaciones accionables & 10\,\% & Discusión, recomendaciones y conclusiones \\
\midrule
Total de la rúbrica oficial & 100\,\% & Cobertura transversal y anexos de trazabilidad \\
\bottomrule
\end{tabularx}}
\apanota{Elaboración basada en \textcite{hun2026terminos}, página~4, secciones 13 y 14, y \textcite{hun2026reto}, página~3. Según la sección~16 de los términos y condiciones (p.~5), los empates se resuelven sucesivamente por optimización y mejora operacional, modelo analítico y rigor metodológico, comprensión del flujo hospitalario y, en último término, valoración conjunta de utilidad, factibilidad y carácter accionable. La redacción en LaTeX, el uso de las normas APA de séptima edición y la extensión del documento obedecen a decisiones editoriales de los autores y no a una imposición de los documentos oficiales del certamen.}
\end{table}

La Tabla~\ref{tab:criterios_oficiales} resume la distribución porcentual de la rúbrica oficial y su correspondencia estructural con los capítulos de este informe. Por rigor normativo, la correspondencia tiene carácter exclusivamente orientador y no introduce subpesos internos ni duplica el peso de un criterio por aparecer tratado de forma transversal entre capítulos. Asimismo, toda distinción entre corridas experimentales archivadas, metadatos inspeccionados del artefacto empaquetado para la web y pruebas locales se mantiene explícita desde la presentación inicial para evitar mezclar evidencia de distinta naturaleza.
